\documentclass[12pt]{article}
\usepackage[utf8]{inputenc}
\usepackage[T1]{fontenc}
\usepackage{hyperref}
\title{Mengaka Language Resources for K-12 Curriculum Development}
\author{Compiled by Claude Sonnet 5.5 (claude-sonnet-5-5) for Andr\'as Kornai}
\date{}
\begin{document}
\maketitle

\section{Introduction}
Mengaka (ISO 639-3 xmg; Glottocode meng1264) is a Bamileke language of the Grassfields (Mbam-Nkam) group in western Cameroon. Other names include Ghap, Megaka, Eghap and Bagam. Speaker estimates conflict: about 20,000 (1993) in one source, 48,000--50,000 in Joshua Project. The language has been written in an indigenous script called Bagam. Little was found in search, and no school materials turned up.

\section{Primary Educational Resources}
No primers, textbooks or school programmes for Mengaka were found in the search results.

\section{Linguistic and Reference Materials}
\begin{itemize}
\item Wikipedia, ``Mengaka language'' (overview, classification, names): \url{https://en.wikipedia.org/wiki/Mengaka_language}
\item Joshua Project, language page: \url{https://joshuaproject.net/index.php/languages/xmg}
\end{itemize}

\section{Religious and Community Resources}
\begin{itemize}
\item Scripture Earth, scripture resources for xmg (Joshua Project reports a New Testament but no complete Bible): \url{https://www.scriptureearth.org/00eng.php?iso=xmg}
\item Joshua Project, Mengaka/Ghap people group in Cameroon: \url{https://joshuaproject.net/people_groups/13672}
\end{itemize}

\section{Digital and Audio Resources}
\begin{itemize}
\item Global Recordings Network, Mengaka audio (including a Bible-stories series): \url{https://globalrecordings.net/en/language/xmg}
\end{itemize}

\section{Curriculum Development Context}
The language is not known to be taught in schools, and no standardised orthography or pedagogical material was found. A curriculum would need new community-led materials, perhaps using the scripture and audio recordings above as starting texts. Local speakers and linguists should be consulted before any development.
\end{document}